\begin{thebibliography}{99}
\bibitem[\protect\citeauthoryear{Lipsitch et al.}{Lipsitch et al.}{2010}]{lipsitch2010}
Lipsitch, M., Tchetgen Tchetgen, E. J., and Cohen, T. 2010. Negative controls: A tool for detecting confounding and bias in observational studies. \textit{Epidemiology} 21(3), 383--388. \url{https://doi.org/10.1097/EDE.0b013e3181d61eeb}.

\bibitem[\protect\citeauthoryear{Schuemie et al.}{Schuemie et al.}{2018}]{schuemie2018}
Schuemie, M. J., Hripcsak, G., Ryan, P. B., Madigan, D., and Suchard, M. A. 2018. Empirical confidence interval calibration for population-level effect estimation studies in observational healthcare data. \textit{PNAS} 115(11), 2571--2577. \url{https://doi.org/10.1073/pnas.1708282114}.

\bibitem[\protect\citeauthoryear{Tchetgen Tchetgen}{Tchetgen Tchetgen}{2014}]{coca2014}
Tchetgen Tchetgen, E. J. 2014. The control outcome calibration approach for causal inference with unobserved confounding. \textit{American Journal of Epidemiology} 179(5), 633--640. \url{https://doi.org/10.1093/aje/kwt303}.

\bibitem[\protect\citeauthoryear{Miao et al.}{Miao et al.}{2024}]{miao2024}
Miao, W., Shi, X., Li, Y., and Tchetgen Tchetgen, E. J. 2024. A confounding bridge approach for double negative control inference on causal effects. \textit{Statistical Theory and Related Fields} 8(4), 262--273. \url{https://doi.org/10.1080/24754269.2024.2390748}.

\bibitem[\protect\citeauthoryear{Park et al.}{Park et al.}{2024}]{park2024}
Park, C., Richardson, D. B., and Tchetgen Tchetgen, E. J. 2024. Single proxy control. \textit{Biometrics} 80(2), ujae027. \url{https://doi.org/10.1093/biomtc/ujae027}.

\bibitem[\protect\citeauthoryear{Hwang et al.}{Hwang et al.}{2022}]{hwang2022}
Hwang, H., Quiroz, J. C., and Gallego, B. 2022. Assessing the effectiveness of empirical calibration under different bias scenarios. \textit{BMC Medical Research Methodology} 22, 208. \url{https://doi.org/10.1186/s12874-022-01687-6}.

\bibitem[\protect\citeauthoryear{Cobzaru et al.}{Cobzaru et al.}{2024}]{cobzaru2024}
Cobzaru, R., Welsch, R., Finkelstein, S., Ng, K., and Shahn, Z. 2024. Bias formulas for violations of proximal identification assumptions in a linear structural equation model. \textit{Journal of Causal Inference} 12(1), 20230039. \url{https://doi.org/10.1515/jci-2023-0039}.

\bibitem[\protect\citeauthoryear{Zheng et al.}{Zheng et al.}{2024}]{zheng2024}
Zheng, J., Wu, J., D'Amour, A., and Franks, A. 2024. Sensitivity to unobserved confounding in studies with factor-structured outcomes. \textit{Journal of the American Statistical Association} 119(547), 2026--2037. \url{https://doi.org/10.1080/01621459.2023.2240053}.

\bibitem[\protect\citeauthoryear{Kuzilek et al.}{Kuzilek et al.}{2017}]{kuzilek2017}
Kuzilek, J., Hlosta, M., and Zdrahal, Z. 2017. Open University Learning Analytics dataset. \textit{Scientific Data} 4, 170171. \url{https://doi.org/10.1038/sdata.2017.171}.

\bibitem[\protect\citeauthoryear{Liu et al.}{Liu et al.}{2023}]{liu2023}
Liu, Z., Liu, Q., Guo, T., Chen, J., Huang, S., Zhao, X., Tang, J., Luo, W., and Weng, J. 2023. XES3G5M: A knowledge tracing benchmark dataset with auxiliary information. \textit{Advances in Neural Information Processing Systems} 36, Datasets and Benchmarks Track. .
\end{thebibliography}
